\documentclass[a4paper,10pt]{article}
\usepackage[utf8]{inputenc} % Corrigido aqui
\usepackage[pt-BR]{babel}
\usepackage[T1]{fontenc}
\usepackage{mathptmx} % Substituto moderno e recomendado para o pacote 'times'
\usepackage{geometry}
\usepackage{graphicx}
\usepackage{amsmath,amssymb}
\usepackage{booktabs}
\usepackage{mdframed}
\usepackage{hyperref}

% Configuração das Margens conforme COBEF 2025
\geometry{
  a4paper,
  top=3cm,
  bottom=2cm,
  left=2cm,
  right=2cm
}

\setlength{\parindent}{0.5cm}
\setlength{\parskip}{0pt}

\begin{document}

% Cabeçalho / Barra Lateral COBEF com Logo UFPR
\begin{mdframed}[topline=false, bottomline=false, rightline=false, linewidth=2.25pt, innerleftmargin=0.3cm, innerrightmargin=0cm]
  \begin{minipage}[c]{0.70\textwidth}
    {\fontsize{14pt}{16pt}\selectfont \bfseries AVALIAÇÃO DA INFLUÊNCIA DO QUEBRA-CAVACO NA FORÇA DE CORTE E ESTUDO DA FORMAÇÃO DO CAVACO NO PROCESSO DE TORNEAMENTO}\\[0.3cm]
    {\small \bfseries Brener Willy Tiago Santos Gomes, Brener.willy@ufpr.br}\\[0.1cm]
    {\small $^{1}$Universidade Federal do Paraná (UFPR), Departamento de Engenharia Mecânica, Curitiba - PR, Brasil}
  \end{minipage}%
  \hfill
  \begin{minipage}[c]{0.25\textwidth}
    \centering
    % Para inserir a logo da UFPR, certifique-se de salvar a imagem como 'logo_ufpr.png' na mesma pasta
    \includegraphics[width=\linewidth]{logo_ufpr.png} 
  \end{minipage}

  \includegraphics{Insertos.png}
  \includegraphics{Nexus-100-2.png}

  \vspace{0.4cm}
  \noindent\textit{\textbf{Resumo:} Este Trabalho foca na compreensão do impacto da força de corte e da geometria de três diferentes quebra-cavacos no processo de torneamento. Após a compreensão das forças de corte resultantes e a morfologia(estudo da forma) dos cavacos gerados, classificados de acordo com a norma ABNT NBR ISO 3685:2017. A instrumentação contou com dinamômetro piezoelétrico . Foi utilizado o torno MAZAK nexus 100-II video da utilização em anexo.nosso objetivo é relacionar o angulo eefetivo de saída com o consumo de porência e a eficiencia do quebra cavaco, para que assim seja possivel fornecer subsídios para a seleção otimizada do ferramental}\\[0.2cm]
  \noindent\textit{\textbf{Palavras-chave:} Usinagem, Quebra-cavaco, Força de corte.}
\end{mdframed}

\vspace{0.5cm}

\section{INTRODUÇÃO}
O processo de torneamento é amplamente utilizado na indústria metalmecânica. A eficiência e a segurança do processo dependem diretamente do controle da cavaco gerado durante o corte. O objetivo deste trabalho é avaliar experimentalmente como a geometria do quebra-cavaco de ferramentas de metal duro influencia a força de corte e a formação do cavaco no torneamento do aço.

\section{REVISÃO BIBLIOGRÁFICA}
O controle do cavaco é essencial em operações automatizadas para evitar severos danos à peça e à ferramenta, além de garantir a segurança do operador. De acordo com Trent e Wright (2000), o quebra-cavaco induz deflexão e tensões de dobramento na fita de cavaco contínua até sua ruptura.

Estudos demonstram que a geometria do quebra-cavaco altera o ângulo efetivo de saída ($\gamma_n$), modificando as componentes da força de corte (Astakhov, 2006).

\section{GEOMETRIA DA FERRAMENTA E CÁLCULO DA FORÇA DE CORTE}

\subsection{Parâmetros Geométricos do Inserto e Suporte}
Para a análise completa do processo de torneamento, determinam-se os seguintes parâmetros geométricos:
\begin{itemize}
    \item \textbf{Ângulo de Posição ($\chi_r$):} Ângulo entre a aresta principal de corte e o plano de avanço. Define a espessura usinada ($h = f \cdot \sin \chi_r$).
    \item \textbf{Ângulo de Folga ($\alpha_n$):} Ângulo entre a superfície de folga da ferramenta e a superfície de corte da peça, evitando o atrito excessivo.
    \item \textbf{Ângulo de Ponta ($\epsilon_r$):} Ângulo entre a aresta principal e a aresta secundária de corte, influenciando a resistência mecânica do inserto:
    \begin{equation}
        \epsilon_r = 180^\circ - (\chi_r + \chi_r')
    \end{equation}
    \item \textbf{Raio de Ponta ($r_s$):} Raio de concordância da ponta do inserto, determinante na rugosidade teórica ($R_z \approx \frac{f^2}{8 r_s}$).
    \item \textbf{Ângulo Efetivo de Saída ($\gamma_n$):} Modificado pelo perfil do quebra-cavaco (rampa ou raio de concordância $r_{qc}$). O ângulo efetivo é calculated por:
    \begin{equation}
        \gamma_{n,efetivo} = \gamma_0 + \arctan\left(\frac{h_{qc}}{w_{qc}}\right)
    \end{equation}
    Onde $h_{qc}$ é a profundidade/altura da rampa do quebra-cavaco e $w_{qc}$ é a largura da canaleta.
\end{itemize}

\subsection{Modelo Teórico da Força Efetiva de Corte ($F_c$)}
A força de corte específica é calculada através da equação clássica de Kienzle:
\begin{equation}
    F_c = k_{c1.1} \cdot b \cdot h^{1-m} \cdot K_{\gamma}
\end{equation}
Substituindo a largura de corte $b = \frac{a_p}{\sin\chi_r}$ e a espessura de corte $h = f \cdot \sin\chi_r$:
\begin{equation}
    F_c = k_{c1.1} \cdot a_p \cdot f^{1-m} \cdot (\sin\chi_r)^{-m} \cdot \left(1 - \frac{\gamma_{efetivo} - \gamma_{ref}}{100}\right)
\end{equation}
Onde $k_{c1.1}$ é a força específica de corte para $b=1\text{ mm}$ e $h=1\text{ mm}$, e $m$ é o expoente característico do material ensaiado.

\section{MATERIAIS E MÉTODOS}
Descreva aqui o procedimento experimental utilizado na aula prática:
\begin{itemize}
    \item \textbf{Material do Corpo de Prova:} Tubo mecânico VMEC 134AP (propriedades mecânicas e composição química conforme norma) adotaod sforça específica de corte 2100N/mm2 eo expoente de Kienzle m ou z = 0,26 logo (1-m)=0,74.
    \item \textbf{Máquina e Instrumentação:} Torno mecânico/CNC MAZAK NEXUS 10-IImunido de dinamômetro piezoelétrico Kistler, modelo Kistler 5070A, ele recebe sinais elétricos(em carga/Coulumbs) e converte esses sinais nume tensão proporcional (Voltdpara medição contínua de forças ($F_c$, $F_f$, $F_p$).
    \item \textbf{Insertos / Quebra-cavacos:} Apresentar a especificação dos 3 insertos testados segundo catálogo do fabricante modelos(WMNA060108 KCR15, WNMG060108 MP KCURO, WMNG060908 FPKCS010).
\end{itemize}

\section{RESULTADOS E DISCUSSÃO}
Foi utilizado o excel para retirar a média de força de corte referênte a cada ferramenta e plotado o gráfico em função do tempo, para cada ferramenta. A seguir, apresento o resultado, nota utilizado o tempo de 50s até 65s onde foi desconsiderado o tempo em vacio da ferramenta, desconsiderado tempo de aproximação e afastamento da ferramenta, onde a força de corte não estabiliza, com exceção do inserto QC_MP e QC_FP, que foi obtido o tempo de 11s até 30s
Apresente aqui os gráficos de forças obtidos do arquivo de dados e classifique as formas dos cavacos segundo a norma **ABNT NBR ISO 3685:2017** (ex.: cavacos em fita, helicoidais, arcos, etc.).

\section{CONCLUSÕES}
Foi verificado que o angulo efetivo de saida proporcionado pela quebra_cavaco reduziu a força de corte sem comprometer a quebra eficiente dos cavacos.
Sintetize os achados verificando se o ângulo efetivo de saída proporcionado pelo quebra-cavaco reduziu a força de corte sem comprometer a quebra eficiente dos cavacos.
A influencia da profundidade de corte ap e do avanço f na força de corte, Kienzel relacionou  o valor da pressão específica de corte varia com a espessura do corte, podendo ser corrigidor em função da variação do angulo de saida

1. FERRARESI, D. Fundamentos da Usinagem dos Metais. Editora Edgard Blücher. 1977 
2.  
DINIZ, A. E. MARCONDES, F. C. COPPINI, N. L. Tecnologia da Usinagem dos Metais. 
Editora Artliber. 5ª edição. 2005 
3. TRENT, E.; M. ;WRIGHT, P. K. Metal Cutting. Butterworth-Heinemann. 2000

\section*{AGRADECIMENTOS}
À Universidade Federal do Paraná (UFPR) e ao Departamento de Engenharia Mecânica pela infraestrutura laboratorial.

\section*{REFERÊNCIAS}
\begin{enumerate}\setlength{\itemsep}{0pt}
    
    \item FERRARESI, D. Fundamentos da Usinagem dos Metais. Editora Edgard Blücher. 1977 
    \item DINIZ, A. E. MARCONDES, F. C. COPPINI, N. L. Tecnologia da Usinagem dos Metais. 
            Editora Artliber. 5ª edição. 2005 
    \item Kienzle, O., 1952. Die Bestimmung von Kräften und Leistungen an Spanenden Werkzeugmaschinen. VDI-Z 94, pp. 299-305.
    \item Trent, E.M. e Wright, P.K., 2000. Metal Cutting. Butterworth-Heinemann, 4th Edition.
    \item Astakhov, V.P., 2006. Tribology of Metal Cutting. Elsevier, London.
    \item ISO 3685, 2017. Tool-life testing with single-point turning tools.
    \item \item ASTM A519 / ISO standards for VMEC 134AP Mechanical Tubing.
\end{enumerate}

\section*{RESPONSABILIDADE PELAS INFORMAÇÕES}
Os autores são os únicos responsáveis pelas informações incluídas neste trabalho.

\end{document}